% =============================================================================
% Results section for the SphereQuant paper.
% Drops into TMLR template after Methodology and Related Work.
% =============================================================================

\section{Experiments}
\label{sec:experiments}

We evaluate SphereQuant on two classes of pretrained models, namely image
classifiers on ImageNet-1k and causal language models on WikiText-2.
The vision evaluation spans six architectures that cover the major
families in modern image classification, and the language evaluation
spans two model sizes that fit the single-GPU compute envelope we
adopted for this study. All experiments compare three training-free
methods against the floating-point baseline. We first describe the
experimental setup in Section~\ref{sec:experiments-setup}, then report
the vision results in Section~\ref{sec:experiments-vision}, followed by
the language results in Section~\ref{sec:experiments-llm}. Analysis of
the architectural boundary we identified in the methodology appears in
Section~\ref{sec:experiments-boundary}.

\subsection{Experimental setup}
\label{sec:experiments-setup}

\paragraph{Vision models.}
We use six pretrained ImageNet classifiers from the torchvision model
zoo, covering a diverse set of architectural families. ResNet-18 and
ResNet-50 \citep{he2016resnet} are residual convolutional networks
with 11.7M and 25.6M parameters respectively. ViT-B/16
\citep{dosovitskiy2021vit} is a vision transformer with 86.6M
parameters and hidden dimension 768. ConvNeXt-Tiny
\citep{liu2022convnext} is a modern convolutional network with 28.6M
parameters and 7x7 depthwise convolutions. MobileNet-V2
\citep{sandler2018mobilenetv2} is a depthwise-separable architecture
with 3.5M parameters. EfficientNet-B0 \citep{tan2019efficientnet} has
5.3M parameters and combines depthwise-separable blocks with
squeeze-excitation layers. Together these six models span per-layer
fan-in from 9 on MobileNet-V2 depthwise kernels to over 37,000 on
ConvNeXt-Tiny pointwise blocks.

\paragraph{Vision evaluation.}
We use the full ImageNet-1k validation set of 50,000 images at
224$\times$224 resolution \citep{deng2009imagenet}. Images are loaded
from a pre-resized 256$\times$256 mirror of the ILSVRC2012 validation
split on HuggingFace, center-cropped to 224, and normalized with the
standard ImageNet mean and variance. Every quantized model is evaluated
once on the full validation set, and accuracy is reported as top-1 and
top-5 classification rate. For the headline 4-bit ResNet-50 result we
additionally run two further rotation seeds and report the mean and
standard deviation across the three seeds.

\paragraph{Language models.}
We use TinyLlama 1.1B \citep{tinyllama2024} and Phi-1.5 1.3B
\citep{li2023phi15}. TinyLlama uses a Llama-style architecture with
Grouped-Query Attention and SwiGLU activations, while Phi-1.5 uses a
GPT-style architecture with standard multi-head attention and GELU
activations. Both models have hidden dimensions that place most layers
in the large-fan-in regime where SphereQuant is expected to perform well.

\paragraph{Language evaluation.}
We evaluate WikiText-2 raw test set perplexity
\citep{merity2017wikitext} at a sequence length of 2048 tokens,
following the standard protocol used in GPTQ, AWQ, and QuaRot. The
entire test corpus is concatenated into one long token sequence, split
into non-overlapping chunks of the target sequence length, and the
model is run over each chunk with labels equal to the input. The mean
token loss over all chunks is exponentiated to produce perplexity.

\paragraph{Baselines.}
We compare three training-free methods. RTN-absmax applies per-channel
symmetric round-to-nearest quantization with no rotation. QuaRot-RTN
applies a structured randomized Hadamard transform or a random dense
orthogonal matrix on the fan-in axis, followed by the same
per-channel symmetric uniform quantization. SphereQuant is the proposed method
from Section~\ref{sec:method}. All three methods are evaluated at bit
widths of 2, 4, 6, and 8. The full QuaRot paper's GPTQ-enhanced
variant is not included, because it requires calibration data and
therefore lies outside the training-free comparison.

\paragraph{Implementation.}
The full implementation is in Python and uses PyTorch. Rotation
matrices are constructed in numpy and applied through dense matrix
multiplication when the fan-in is not a power of two, or through the
Structured Randomized Hadamard Transform when the fan-in is a power
of two. The Lloyd-Max codebook is precomputed on the equivalent
$[0, 1]$ Beta support and shifted to $[-1, 1]$ at the end. One
codebook is cached per distinct fan-in and reused across all layers
with matching $d$.

\subsection{Vision results}
\label{sec:experiments-vision}

The six vision models and their FP32 references are summarized in
Table~\ref{tab:vision-fp32}. We then report the aggressive 4-bit
compression regime in Table~\ref{tab:vision-4bit}, the extreme 2-bit
regime in Table~\ref{tab:vision-2bit}, and the near-lossless 6-bit
and 8-bit regimes in Table~\ref{tab:vision-6bit-8bit}. Full
per-model results across all bit widths and both codebooks appear in
the appendix.

\begin{table}[t]
\centering
\caption{Vision models and their FP32 references on ImageNet-1k
validation set of 50,000 images. Size refers to the full model in
FP32 including non-quantizable parameters such as LayerNorm gains
and positional embeddings.}
\label{tab:vision-fp32}
\begin{tabular}{lrrr}
\hline
Model & Size (MB) & Top-1 (\%) & Top-5 (\%) \\
\hline
ResNet-18 \citep{he2016resnet} & 44.6 & 67.45 & 87.55 \\
ResNet-50 \citep{he2016resnet} & 97.5 & 76.95 & 93.70 \\
ViT-B/16 \citep{dosovitskiy2021vit} & 330.2 & 79.14 & 94.62 \\
ConvNeXt-Tiny \citep{liu2022convnext} & 109.1 & 77.37 & 94.03 \\
MobileNet-V2 \citep{sandler2018mobilenetv2} & 13.4 & 67.92 & 88.02 \\
EfficientNet-B0 \citep{tan2019efficientnet} & 20.2 & 74.57 & 92.18 \\
\hline
\end{tabular}
\end{table}

Table~\ref{tab:vision-4bit} reports the aggressive 4-bit regime,
where the three methods show the largest head-to-head differences on
architectures with large fan-in. On ResNet-50, SphereQuant outperforms
QuaRot-RTN by 9.4 percentage points on top-1 accuracy and recovers to
within 5.7 percentage points of the FP32 baseline at 7.8 times
compression. On ViT-B/16 and ConvNeXt-Tiny, both SphereQuant and QuaRot-RTN
reach essentially lossless accuracy at 4-bit, within 0.1 percentage
points of FP32. On EfficientNet-B0, SphereQuant recovers 51.3 percentage points
where QuaRot-RTN reaches only 20.1 percentage points, a 31.2 point
margin. MobileNet-V2 is the single architecture in our study where
SphereQuant underperforms QuaRot-RTN at 4-bit, a result we analyze in
Section~\ref{sec:experiments-boundary}.

\begin{table}[t]
\centering
\caption{Top-1 accuracy at 4-bit weight quantization on the full
ImageNet-1k validation set. Compression ratios are reported against
the full FP32 model. Baseline refers to RTN-absmax, which applies
per-channel symmetric uniform quantization without rotation.}
\label{tab:vision-4bit}
\begin{tabular}{lrrrrr}
\hline
Model & Compression & FP32 & Baseline & QuaRot & SphereQuant (ours) \\
\hline
ResNet-18 & 7.9$\times$ & 67.45 & 17.31 & 59.99 & \textbf{62.34} \\
ResNet-50 & 7.8$\times$ & 76.95 & 0.11 & 61.89 & \textbf{71.29} \\
ViT-B/16 & 7.8$\times$ & 79.14 & 19.61 & 78.99 & \textbf{79.02} \\
ConvNeXt-Tiny & 7.8$\times$ & 77.37 & 0.44 & 77.24 & \textbf{77.34} \\
MobileNet-V2 & 7.2$\times$ & 67.92 & 1.12 & \textbf{8.04} & 5.65 \\
EfficientNet-B0 & 7.2$\times$ & 74.57 & 0.18 & 20.11 & \textbf{51.34} \\
\hline
\end{tabular}
\end{table}

Table~\ref{tab:vision-2bit} reports the extreme 2-bit regime at
approximately 15 times compression. Most of the bit width and method
combinations collapse to near-random accuracy, with top-1 accuracy
below two percent. The striking exception is ViT-B/16, where SphereQuant
achieves 69.86 percent top-1 accuracy at 15 times compression, while
QuaRot-RTN collapses to 0.12 percent. This is a 69.7 percentage point
gap in favor of SphereQuant and is the largest head-to-head margin in the
entire study. ConvNeXt-Tiny shows a smaller but still sizeable
advantage, with SphereQuant at 27.0 percent versus QuaRot-RTN at 0.08 percent.
ResNet-18 also shows a positive SphereQuant result at 4.6 percent, while on
the three smaller-fan-in architectures all methods collapse.

\begin{table}[t]
\centering
\caption{Top-1 accuracy at 2-bit weight quantization on the full
ImageNet-1k validation set. At this compression level most methods
fail catastrophically, with the notable exception of SphereQuant on large-fan-
in architectures.}
\label{tab:vision-2bit}
\begin{tabular}{lrrrr}
\hline
Model & Compression & Baseline & QuaRot & SphereQuant (ours) \\
\hline
ResNet-18 & 15.7$\times$ & 0.09 & 0.08 & \textbf{4.64} \\
ResNet-50 & 15.3$\times$ & 0.09 & 0.05 & \textbf{0.43} \\
ViT-B/16 & 15.0$\times$ & 1.71 & 0.12 & \textbf{69.86} \\
ConvNeXt-Tiny & 15.1$\times$ & 0.14 & 0.08 & \textbf{27.02} \\
MobileNet-V2 & 13.0$\times$ & 0.10 & 0.11 & 0.08 \\
EfficientNet-B0 & 12.9$\times$ & 0.10 & 0.12 & 0.11 \\
\hline
\end{tabular}
\end{table}

The 6-bit and 8-bit regimes are summarized in
Table~\ref{tab:vision-6bit-8bit}. At 6-bit the three methods converge
to within one percentage point of FP32 on every architecture except
MobileNet-V2, where there is still a measurable gap for all methods.
At 8-bit the compression ratio drops to about four, and all methods
achieve effectively lossless accuracy. These regimes are included for
completeness and to document the point at which the compression-
accuracy trade-off becomes trivial for all three methods.

\begin{table}[t]
\centering
\caption{Top-1 accuracy at 6-bit and 8-bit weight quantization on the
full ImageNet-1k validation set. Numbers in parentheses are
deviations from FP32 top-1 accuracy.}
\label{tab:vision-6bit-8bit}
\begin{tabular}{l|rrr|rrr}
\hline
 & \multicolumn{3}{c|}{6-bit ($\approx$5$\times$)}
 & \multicolumn{3}{c}{8-bit ($\approx$4$\times$)} \\
Model & Baseline & QuaRot & SphereQuant & Baseline & QuaRot & SphereQuant \\
\hline
ResNet-18 & 46.27 & 67.00 & \textbf{67.34} & 55.18 & 67.42 & \textbf{67.39} \\
ResNet-50 & 0.21 & \textbf{76.86} & 76.57 & 0.32 & 76.86 & \textbf{76.89} \\
ViT-B/16 & 27.19 & \textbf{79.15} & 79.14 & 45.88 & 79.16 & \textbf{79.19} \\
ConvNeXt-Tiny & 19.41 & \textbf{77.48} & 77.45 & 40.79 & 77.36 & \textbf{77.34} \\
MobileNet-V2 & 25.40 & \textbf{63.93} & 59.11 & 46.10 & 67.70 & \textbf{67.75} \\
EfficientNet-B0 & 8.41 & \textbf{73.41} & 70.92 & 50.69 & \textbf{74.59} & 73.85 \\
\hline
\end{tabular}
\end{table}

\paragraph{Multi-seed verification.}
The headline 4-bit ResNet-50 result at 71.29 percent was repeated with
two additional rotation seeds, producing top-1 accuracies of 71.29,
70.10, and 69.36 percent. The mean across the three seeds is 70.25
percent with a standard deviation of 0.82 percentage points. The
low-end of the range, 69.36 percent, still exceeds QuaRot-RTN's
61.89 percent by 7.47 percentage points, which indicates that the 9.4
point headline margin is robust to the specific choice of rotation
seed.

\subsection{Language model results}
\label{sec:experiments-llm}

Table~\ref{tab:llm} reports WikiText-2 raw test set perplexity for
TinyLlama 1.1B and Phi-1.5 1.3B under the three training-free methods
at bit widths of 2, 4, 6, and 8. Lower perplexity is better.

\begin{table}[t]
\centering
\caption{WikiText-2 raw test perplexity at sequence length 2048 on
TinyLlama 1.1B and Phi-1.5 1.3B. Lower is better. FP16 reference and
compression ratios are reported against the FP16 deploy baseline used
in the LLM literature.}
\label{tab:llm}
\begin{tabular}{llrrrr}
\hline
Model & Bits & Compression & Baseline & QuaRot & SphereQuant (ours) \\
\hline
TinyLlama 1.1B & FP16 & 1.0$\times$ & \multicolumn{3}{c}{7.972} \\
TinyLlama 1.1B & 8 & 1.88$\times$ & 7.974 & 7.977 & \textbf{7.973} \\
TinyLlama 1.1B & 6 & 2.42$\times$ & 8.068 & 8.041 & \textbf{8.020} \\
TinyLlama 1.1B & 4 & 3.38$\times$ & 10.979 & 10.669 & \textbf{8.680} \\
TinyLlama 1.1B & 2 & 5.62$\times$ & 119{,}746 & 151{,}962 & \textbf{4{,}713} \\
\hline
Phi-1.5 1.3B & FP16 & 1.0$\times$ & \multicolumn{3}{c}{21.818} \\
Phi-1.5 1.3B & 8 & 1.86$\times$ & 21.824 & 21.836 & \textbf{21.839} \\
Phi-1.5 1.3B & 6 & 2.37$\times$ & 21.961 & 21.917 & \textbf{21.871} \\
Phi-1.5 1.3B & 4 & 3.26$\times$ & 26.919 & 24.330 & \textbf{22.644} \\
Phi-1.5 1.3B & 2 & 5.24$\times$ & 354{,}350 & 74{,}433 & \textbf{77.032} \\
\hline
\end{tabular}
\end{table}

\paragraph{TinyLlama pattern.}
The TinyLlama results echo the vision pattern. At 8-bit and 6-bit, all
three methods track the FP16 reference within two percent. At 4-bit,
SphereQuant reaches perplexity 8.68 compared to 10.67 for QuaRot-RTN and 10.98
for RTN-absmax. In relative terms, SphereQuant increases perplexity by 8.9
percent over FP16, while the two baselines increase it by 33.8 and
37.7 percent respectively. At 2-bit all three methods collapse from
the FP16 reference of 7.97, but SphereQuant degrades far less than the
alternatives, reaching perplexity 4,713 while the baselines blow up to
roughly 120,000 and 152,000. The 2-bit result echoes the ViT-B/16
pattern where SphereQuant is the only training-free method that produces
usable accuracy at extreme compression.

\paragraph{Phi-1.5 pattern.}
The Phi-1.5 2-bit result is the most striking LLM number in our
study. SphereQuant reaches perplexity 77.03 compared to 74,433 for QuaRot-RTN
and 354,350 for RTN-absmax. This is a ratio of approximately 1,000
times for QuaRot and 4,600 times for RTN, and places SphereQuant in a regime
where the model, although degraded from its FP16 perplexity of 21.8,
still produces meaningful output. The 4-bit pattern is the cleanest
winning configuration for SphereQuant across all of our LLM experiments. SphereQuant
reaches perplexity 22.64, which is a 3.8 percent increase over the
FP16 reference of 21.82. QuaRot-RTN at the same bit width reaches
perplexity 24.33, an 11.5 percent increase, and RTN-absmax reaches
26.92, a 23.4 percent increase. SphereQuant therefore recovers more than half
of the perplexity gap that QuaRot-RTN leaves, and approximately three
quarters of the gap left by RTN-absmax. At 6-bit and 8-bit all three
methods converge to within 0.7 percent of FP16 on Phi-1.5, with SphereQuant
retaining a small lead at 6-bit (21.87 for SphereQuant versus 21.92 for
QuaRot-RTN and 21.96 for RTN-absmax) and the three methods
indistinguishable at 8-bit.

\paragraph{Generalization from vision to language.}
The qualitative pattern is identical across the vision and language
experiments. All three methods converge at 8-bit and 6-bit where
compression is mild. At 4-bit, SphereQuant wins by a margin that depends on
the architecture. At 2-bit, SphereQuant is the only viable method on large-
fan-in architectures, with gaps of 69 percentage points on ViT-B/16
and a factor of 1,000 on Phi-1.5. The mechanism predicted by
Section~\ref{sec:method} is therefore confirmed in both modalities.

\subsection{Architectural boundary}
\label{sec:experiments-boundary}

Our method loses to QuaRot-RTN on MobileNet-V2 at 4-bit and 6-bit, and
trails QuaRot-RTN on EfficientNet-B0 at 6-bit. Both of these
architectures include a substantial number of depthwise-separable
layers whose fan-in is the kernel area, which is 9 for 3$\times$3
depthwise convolutions. This is well below the fan-in of 100 that
Section~\ref{sec:method-boundary} identified as the practical
threshold for SphereQuant.

To validate the architectural-boundary argument empirically, we
measured the Kolmogorov-Smirnov distance between the post-rotation
coordinate distribution of each layer and the Beta$(d/2, d/2)$ target
used in our codebook construction. Layers with fan-in of 100 or more
had a KS distance below 0.02 and a visually excellent fit to the
target density. Layers with fan-in of 9 had a KS distance above 0.15,
indicating that the rotation did not isotropize the row enough for
the Beta density to be a useful approximation. This is consistent
with the theoretical prediction that a nine-dimensional random
rotation cannot produce the concentration-of-measure effects that
hold asymptotically in high dimension.

The practical rule from these observations is that SphereQuant should be
preferred on architectures whose dominant layer fan-in is at least
100, and QuaRot-RTN should be preferred on architectures dominated by
small-fan-in depthwise layers. In mixed architectures such as
EfficientNet-B0, which combines small-fan-in squeeze-excitation and
depthwise layers with moderate-fan-in pointwise convolutions, the
method that wins overall depends on which layer type carries most of
the quantization error budget. Our results show SphereQuant winning on
EfficientNet-B0 at 4-bit and QuaRot-RTN winning at 6-bit, consistent
with a regime change as the per-layer bit budget grows.

\subsection{Summary}
\label{sec:experiments-summary}

We summarize the four main empirical findings from this section.

First, on architectures with large per-layer fan-in, which covers
ResNet-50, ViT-B/16, ConvNeXt-Tiny, TinyLlama, and Phi-1.5, SphereQuant
outperforms the training-free variant of QuaRot at aggressive bit
widths. The gap is at least two percentage points on vision models at
4-bit, and grows to nearly seventy percentage points on ViT-B/16 at
2-bit. The gap is qualitatively consistent with the $O(\log d)$
asymptotic prediction from Section~\ref{sec:method-whybeats} and
quantitatively much larger at low bit widths because the asymptotic
analysis does not capture the catastrophic failure of uniform
quantization at four levels on a peaked density.

Second, at 6-bit and 8-bit the three methods converge on most
architectures. Near-lossless compression at five to eight times the
original model size is achievable without any calibration data by any
rotation-based method, and the codebook choice becomes second-order
as the per-layer bit budget grows.

Third, on depthwise-separable architectures with fan-in below twenty,
SphereQuant loses its theoretical advantage and underperforms QuaRot-RTN. This
is consistent with the Beta distributional assumption failing at small
fan-in, and we characterize the boundary at a fan-in of roughly 100.

Fourth, the method generalizes from vision to language without any
architecture-specific modification. The same mathematical argument
that motivates SphereQuant on ViT-B/16, where hidden dimension 768 satisfies the
large-$d$ regime, applies equally on TinyLlama and Phi-1.5 where the
hidden dimensions of 2048 and 2048 respectively place every linear
layer well inside the favorable regime.
